% ============================================================================
\chapter{Communicuer avec le serveur : fetch, SSE, WebSockets, workers}
\label{ch:js-reseau}
% ============================================================================

\begin{objectif}
\begin{itemize}
  \item Ma\^itriser \texttt{fetch} dans tous ses usages~: JSON, en-t\^etes, erreurs, annulation,
        progression, t\'el\'echargement, upload.
  \item Structurer une couche API testable et agr\'eable (client HTTP maison, intercepteurs, cache).
  \item Consommer des flux temps r\'eel~: Server-Sent Events, WebSockets, reconnexion.
  \item D\'eplacer le travail lourd hors du fil principal~: Web Workers, OffscreenCanvas, WASM.
  \item Comprendre CORS et le r\'egler c\^ot\'e client \emph{et} c\^ot\'e serveur.
\end{itemize}
\end{objectif}

\section{Le strict minimum bien fait}

\begin{codejs}
async function getJSON(url, options = {}) {
  const ctrl = new AbortController();
  const minuterie = setTimeout(() => ctrl.abort(new Error("Timeout 10s")), 10_000);
  try {
    const res = await fetch(url, {
      signal: ctrl.signal,
      credentials: "same-origin",        // envoyer les cookies de session du meme site
      headers: { accept: "application/json", ...options.headers },
      ...options,
    });
    // fetch ne rejette PAS sur 404/500 : c'est a toi de le faire
    if (!res.ok) {
      const corps = await res.text().catch(() => "");
      throw new Error(`HTTP ${res.status} ${res.statusText} sur ${url}${corps ? ` -- ${corps.slice(0, 200)}` : ""}`);
    }
    if (res.status === 204) return null;
    const type = res.headers.get("content-type") ?? "";
    return type.includes("json") ? res.json() : res.text();
  } finally {
    clearTimeout(minuterie);
  }
}

const creer = (donnees) => getJSON("/api/articles", {
  method: "POST",
  headers: { "content-type": "application/json" },
  body: JSON.stringify(donnees),
});
const supprimer = (id) => getJSON(`/api/articles/${id}`, { method: "DELETE" });
\end{codejs}

\begin{tabz}{lX}
  \headrow \thd{Besoin} \& \thd{Solution} \\
  \midrule
  Envoi d'un formulaire simple \& \verb|await fetch(url, { method: "POST", body: new FormData(form) })| \\
  Upload d'un fichier \& \verb|body: formData| avec \verb|File| dedans (jamais \verb|Content-Type| manuel) \\
  Progression d'envoi \& \verb|XMLHttpRequest.upload.onprogress| ou \verb|fetch| + stream request (\verb|duplex: "half"|) \\
  Progression de r\'eception \& \verb|res.body| + compteur d'octets (stream) \\
  T\'el\'echargement binaire \& \verb|const blob = await res.blob()| puis \verb|URL.createObjectURL| \\
  Requ\^ete annulable \& \verb|AbortController| (aussi pour \'eviter les r\'eponses hors ordre) \\
  Cache HTTP fin \& \verb|cache: "no-store"|, \verb|"force-cache"|, \verb|"reload"| \\
  Authentification par jeton \& en-t\^ete \verb|authorization: Bearer ...| \\
  Cookies de session \& \verb|credentials: "include"| + \verb|SameSite=Lax| c\^ot\'e serveur \\
  Streaming de texte (IA) \& \verb|res.body.pipeThrough(new TextDecoderStream())| + \verb|for await| \\
\end{tabz}

\subsection{Un mini-client HTTP r\'eutilisable}

\begin{codejs}
// api.js --- le fichier que j'ecris en premier dans tout projet
const BASE = import.meta.env?.VITE_API_URL ?? "/api";

class ErreurAPI extends Error {
  constructor(message, { status, code, details } = {}) {
    super(message);
    this.name = "ErreurAPI"; this.status = status; this.code = code; this.details = details ?? null;
  }
}

function creerClient({ baseURL = BASE, getJeton, on401 } = {}) {
  const controles = new AbortController();
  async function requete(methode, chemin, { corps, signaux, enTetes, raw = false, idempotent = methode === "GET" } = {}) {
    let tentative = 0;
    while (true) {
      tentative++;
      const url = new URL(chemin.replace(/^\//, ""), baseURL.endsWith("/") ? baseURL : `${baseURL}/`).href;
      const jeton = await getJeton?.();
      let res;
      try {
        res = await fetch(url, {
          method: methode,
          signal: AbortSignal.any([...(signaux ?? []), controles.signal]),
          headers: {
            accept: "application/json",
            ...(corps !== undefined && !(corps instanceof FormData) ? { "content-type": "application/json" } : {}),
            ...(jeton ? { authorization: `Bearer ${jeton}` } : {}),
            ...enTetes,
          },
          ...(corps !== undefined ? { body: corps instanceof FormData || corps instanceof Blob ? corps : JSON.stringify(corps) } : {}),
        });
      } catch (err) {
        if (err.name === "AbortError") throw new DOMException("Requete annulee", "AbortError");
        if (idempotent && tentative < 3) { await new Promise((r) => setTimeout(r, 200 * 2 ** tentative)); continue; }
        throw new ErreurAPI("Serveur injoignable", { code: "reseau", details: err.message });
      }
      if (res.status === 401) { await on401?.(); throw new ErreurAPI("Session expiree", { status: 401 }); }
      if (!res.ok) {
        const payload = await res.json().catch(() => ({}));
        throw new ErreurAPI(payload.message ?? `HTTP ${res.status}`, { status: res.status, code: payload.code, details: payload.errors });
      }
      if (raw) return res;
      if (res.status === 204) return null;
      return (res.headers.get("content-type") ?? "").includes("json") ? res.json() : res.text();
    }
  }
  return {
    get: (c, o) => requete("GET", c, o),
    post: (c, corps, o) => requete("POST", c, { ...o, corps }),
    patch: (c, corps, o) => requete("PATCH", c, { ...o, corps }),
    put: (c, corps, o) => requete("PUT", c, { ...o, corps }),
    delete: (c, o) => requete("DELETE", c, o),
    brut: (methode, c, o) => requete(methode, c, { ...o, raw: true }),
    annulerTout: () => controles.abort(),
  };
}

export const api = creerClient({
  getJeton: () => sessionStorage.getItem("token"),
  on401: async () => { sessionStorage.clear(); location.assign("/connexion"); },
});
const articles = await api.get("/articles?page=1&tri=recent");
\end{codejs}

\begin{note}
Ne mets jamais un secret (cl\'e d'API priv\'ee, mot de passe base de donn\'ees) dans un bundle
JavaScript~: le navigateur le montre \`a quiconque ouvre l'onglet Sources. Tout ce qui doit
rester secret passe par ton propre serveur.
\end{note}

\section{CORS~: l'erreur que tout le monde rencontre}

Le navigateur refuse qu'une page de \verb|https://a.dev| lise une r\'eponse de
\verb|https://b.dev| sans autorisation explicite du serveur. Ce n'est \emph{pas} une s\'ecurit\'e
pour le serveur (il re\c{c}oit bien la requ\^ete)~: c'est une protection du \emph{navigateur}
pour l'utilisateur.

\begin{codejs}
// Cote serveur (ici en Express) --- les en-tetes qui debloquent la situation
app.use((req, res, next) => {
  const autorises = new Set(["http://localhost:5173", "https://nour.dev"]);
  const origine = req.headers.origin;
  if (autorises.has(origine)) {
    res.setHeader("access-control-allow-origin", origine);
    res.setHeader("vary", "origin");
    res.setHeader("access-control-allow-credentials", "true");
    res.setHeader("access-control-allow-methods", "GET,POST,PUT,PATCH,DELETE,OPTIONS");
    res.setHeader("access-control-allow-headers", "content-type,authorization,x-csrf-token");
    res.setHeader("access-control-max-age", "600");
  }
  if (req.method === "OPTIONS") return res.sendStatus(204);   // pre-vol
  next();
});
\end{codejs}

\etape{Lire le vrai message}
Dans la console~: \verb|Access to fetch at '...' has been blocked by CORS policy: No
'Access-Control-Allow-Origin' header...|. La cause est dans la \emph{deuxi\`eme} ligne~:
soit l'en-t\^ete manque, soit \verb|credentials: "include"| est utilis\'e avec
\verb|allow-origin: *| (interdit), soit un en-t\^ete personnalis\'e n'est pas list\'e dans
\verb|allow-headers|.

\etape{Pendant le d\'eveloppement, la vraie solution}
Ne touche pas au serveur~: fais passer les appels par le serveur de d\'eveloppement
(m\^eme origine, donc pas de CORS).

\begin{codejson}
{
  "server": {
    "proxy": {
      "/api": { "target": "http://localhost:8000", "changeOrigin": true }
    }
  },
  "scripts": { "dev": "vite", "dev:api": "php -S localhost:8000 -t public" }
}
\end{codejson}

\begin{codeplain}
// vite.config.js --- la meme idee avec Vite
import { defineConfig } from "vite";
export default defineConfig({
  server: {
    proxy: {
      "/api": { target: "http://localhost:8000", changeOrigin: true },
      "/ws": { target: "ws://localhost:8000", ws: true },
    },
  },
});
\end{codeplain}

\section{Server-Sent Events~: le temps r\'eel le plus simple du monde}

Une connexion HTTP classique, un flux texte infini, reconnexion automatique par le navigateur.
Parfait pour des notifications~; insuffisant pour du bidirectionnel.

\begin{codejs}
const sse = new EventSource("/api/stream");
sse.onmessage = (e) => rendre(e.data);                 // evenement sans nom
sse.addEventListener("notification", (e) => {
  const data = JSON.parse(e.data);
  notifier(data);
});
sse.onerror = () => { /* le navigateur se reconnecte seul ; a toi de gerer l'etat "hors ligne" */ };
// Cote serveur (Node/Express) :
// res.writeHead(200, { "content-type": "text/event-stream", "cache-control": "no-cache", connection: "keep-alive" });
// res.write(`event: notification\ndata: ${JSON.stringify(objet)}\n\n`);
// res.write(": ping\n\n");     // commentaire de maintien en vie toutes les 20-30 s
\end{codejs}

\section{WebSockets~: bidirectionnel}

\begin{codejs}
class Socket {
  #ws; #file = []; #retry = 0; #ctrl;
  constructor(url, { surMessage, surOuverture } = {}) {
    Object.assign(this, { url, surMessage, surOuverture });
  }
  ouvrir() {
    this.#ws = new WebSocket(this.url);
    this.#ws.onopen = () => {
      this.#retry = 0;
      this.#file.splice(0).forEach((m) => this.#ws.send(m));
      this.surOuverture?.();
    };
    this.#ws.onmessage = (e) => {
      const brut = typeof e.data === "string" ? e.data : lireBlob(e.data);
      try { this.surMessage?.(JSON.parse(brut)); }
      catch { this.surMessage?.(brut); }
    };
    this.#ws.onclose = (e) => {
      if (e.code === 1000) return;                       // fermeture propre
      this.#retry++;
      const attente = Math.min(30_000, 500 * 2 ** this.#retry);
      console.warn(`WebSocket ferme (${e.code}), reconnexion dans ${attente} ms`);
      this.#t = setTimeout(() => this.ouvrir(), attente + Math.random() * 250);
    };
    this.#ws.onerror = () => this.#ws.close();
  }
  envoyer(objet) {
    const msg = typeof objet === "string" ? objet : JSON.stringify(objet);
    if (this.#ws?.readyState === WebSocket.OPEN) this.#ws.send(msg);
    else this.#file.push(msg);                            // mettre en attente, pas jeter
  }
  fermer() { clearTimeout(this.#t); this.#ws?.close(1000); }
  get enLigne() { return this.#ws?.readyState === WebSocket.OPEN; }
}
const lireBlob = async (blob) => blob.text();   // ou arrayBuffer() pour du binaire
const ws = new Socket("wss://nour.dev/ws", { surMessage: (m) => console.log("recu", m) });
ws.envoyer({ type: "abonner", salon: "general" });
navigator.addEventListener("online", () => !ws.enLigne && ws.ouvrir());
\end{codejs}

\begin{note}
WebSocket ne propose ni authentification par cookie de la m\^eme fa\c{c}on ni reconnexion~: tu les
codes. C'est pourquoi, pour du simple temps r\'eel, SSE (serveur \$\rightarrow\$ client) ou une
solution h\'eberg\'ee (\link{Pusher}{https://pusher.com/}, \link{Ably}{https://www.ably.io/},
\link{Supabase Realtime}{https://supabase.com/}) est souvent le choix raisonnable.
C\^ot\'e serveur Node, la biblioth\`eque de r\'ef\'erence est
\link{ws}{https://github.com/websockets/ws} ou \link{Socket.IO}{https://socket.io/}.
\end{note}

\section{Web Workers : le CPU en parall\`ele}
\label{sec:workers}

Un worker est un fil JavaScript s\'epar\'e~: pas d'acc\`es au DOM, mais un CPU entier. On \'echange
des messages (structur\'es, clonables) et des \emph{transferables} ( ArrayBuffer, sans copie).

\begin{codeplain}
// poids-lourd.js (le fichier du worker)
self.onmessage = async (e) => {
  const { type, payload } = e.data;
  if (type === "ligner") {
    const t0 = performance.now();
    const resultat = calcul(payload);
    self.postMessage({ type: "resultat", payload: resultat, ms: performance.now() - t0 });
  }
};
function calcul(n) { let s = 0; for (let i = 0; i < n; i++) s += Math.sqrt(i); return s; }

// main.js
const w = new Worker(new URL("./poids-lourd.js", import.meta.url), { type: "module" });
w.postMessage({ type: "ligner", payload: 5e7 });
w.onmessage = (e) => { document.querySelector("#out").textContent = `${e.data.payload} en ${e.data.ms.toFixed(0)} ms`; };
// w.terminate();   // tuer le worker
\end{codeplain}

\begin{itemize}
  \item \link{Comlink}{https://github.com/GoogleChromeLabs/comlink} rend l'appel de worker
        presque banal (promesses \`a travers les fronti\`eres).
  \item \verb|navigator.hardwareConcurrency| donne le nombre de coeurs~; un
        \verb|WorkerPool| de \$n-1\$ workers traite les gros jeux de donn\'ees en parall\`ele.
  \item \textbf{SharedArrayBuffer} + \verb|Atomics| partagent r\'eellement la m\'emoire
        (n\'ecessite les en-t\^etes \verb|Cross-Origin-Opener-Policy| / \verb|Embedder-Opener-Policy|).
  \item \textbf{OffscreenCanvas} d\'el\`egue le rendu WebGL/canvas \`a un worker~: l'interface ne
        rame plus pendant les animations.
  \item \textbf{OPFS} (\link{Origin Private File System}{https://developer.mozilla.org/en-US/docs/Web/API/File\_System\_API/Origin\_private\_file\_system})
        est un stockage de fichiers rapide, isol\'e, utilisable depuis un worker~: c'est la voie
        s\'erieuse pour des donn\'ees volumineuses c\^ot\'e client.
  \item WebAssembly~: le langage des autres langages (Rust, C, Go). Appeler du WASM depuis JS
        est triviale~; l'\'ecrire est un autre livre. D\'emo rapide~:
        \link{Rust et wasm-pack}{https://rustwasm.github.io/wasm-pack/}.
\end{itemize}

\begin{exercice}
\begin{enumerate}
  \item Consomme une API publique r\'eelle --- \url{https://api.github.com/repos/vitejs/vite} ou
        \url{https://api.open-meteo.com/v1/forecast?latitude=14.79\&longitude=-16.67\&current=temperature\\\_2m}
        avec le mini-client~: affiche la temp\'erature de Thi\`es, g\`ere erreur r\'eseau et hors ligne.
  \item Ajoute un cache~: garde la derni\`ere r\'eponse dans \verb|sessionStorage| avec son \^age,
        affiche-la imm\'ediatement, puis rafra\^ichis en arri\`ere-plan (strat\'egie
        \og ~stale-while-revalidate~ \fg{}).
  \item Impl\'emente l'annulation~: taper une nouvelle recherche annule la requ\^ete pr\'ec\'edente,
        et une r\'eponse hors ordre est jet\'ee.
  \item \'Ecris un affichage temps r\'eel~: un petit serveur Node SSE qui \'emet l'heure toutes les
        secondes, un client qui affiche, d\'etecte la d\'econnexion et affiche \og ~hors ligne~ \fg{}.
  \item D\'eplace un calcul lourd (somme de 10\textsuperscript{8} racines carr\'ees, ou tri de
        10\textsuperscript{6} nombres) dans un worker et mesure le gain sur le FPS d'une
        animation~: utilise le panneau \clic{Performance}.
  \item Ajoute la reprise sur erreur r\'eseau~: file d'attente des \'ecritures locales, renvoi
        automatique quand \verb|online| revient, affichage du nombre d'op\'erations en attente.
\end{enumerate}
\end{exercice}

\begin{ressources}
\begin{itemize}
  \item MDN --- fetch --- \url{https://developer.mozilla.org/fr/docs/Web/API/Fetch\\\_API/Using\\\_Fetch}
  \item MDN --- CORS --- \url{https://developer.mozilla.org/fr/docs/Web/HTTP/Guides/CORS}
  \item \link{Async API tutorial (Node, SSE, WS)}{https://github.com/nodejs/node} --- docs \og ~http~ \fg{}
  \item \link{web.dev : patterns de reseau}{https://web.dev/patterns/}
  \item \link{MDN : Using Web Workers}{https://developer.mozilla.org/fr/docs/Web/API/Web\_Workers\_API/Using\_web\_workers}
\end{itemize}
\end{ressources}

\index{fetch}\index{CORS}\index{WebSocket}\index{worker}
